\section{Zadanie 6 (Mateusz Wojaczek)}
Pokaż, że jeżeli zdarzenia $A_1, A_2, \ldots, A_n$ są niezależne, 
to zdarzenia 
\[
\Omega \setminus A_1,\, \Omega \setminus A_2,\, \ldots,\, \Omega \setminus A_n
\]
również są niezależne.

\subsection*{Rozwiązanie}

\subsection*{Założenie (niezależność zdarzeń $A_1, \dots, A_n$)}

\[
\forall I \subseteq \{1,2,\dots,n\}, \quad
P\!\left(\bigcap_{i \in I} A_i \right) = \prod_{i \in I} P(A_i).
\]

\subsection*{Teza do udowodnienia}

Dla dowolnego $I \subseteq \{1, \dots, n\}$ zachodzi
\[
P\!\left(\bigcap_{i \in I} (\Omega \setminus A_i)\right)
= \prod_{i \in I} P(\Omega \setminus A_i).
\]

\subsection*{Dowód}

Niech $I \subseteq [n]$. Rozpoczynamy od przekształcenia przecięcia dopełnień przy użyciu prawa De Morgana:

\begin{align*}
P\!\left( \bigcap_{i \in I} (\Omega \setminus A_i) \right)
&= P\!\left( \Omega \setminus \bigcup_{i \in I} A_i \right) 
\text{(De Morgan)} \\
&= P(\Omega) - P\!\left( \bigcup_{i \in I} A_i \right) 
= 1 - P\!\left( \bigcup_{i \in I} A_i \right).
\end{align*}


\[
P\!\left( \bigcap_{i \in I} (\Omega \setminus A_i) \right)
= 1 - \sum_{\substack{J \subseteq I \\ J \neq \emptyset}} (-1)^{|J|+1} P\!\left( \bigcap_{i \in J} A_i \right).
\]

Ponieważ zdarzenia $A_i$ są niezależne

\[
P\!\left( \bigcap_{i \in J} A_i \right) = \prod_{i \in J} P(A_i),
\]

więc

\[
P\!\left( \bigcap_{i \in I} (\Omega \setminus A_i) \right)
= 1 - \sum_{\substack{J \subseteq I \\ J \neq \emptyset}} (-1)^{|J|+1} \prod_{i \in J} P(A_i).
\]

Teraz zauważamy, że suma odpowiada funkcji tworzącej iloczynu:

\[
\sum_{\substack{J \subseteq I \\ J \neq \emptyset}} (-1)^{|J|+1} \prod_{i \in J} P(A_i)
= 1 - \prod_{i \in I} \big( 1 - P(A_i) \big).
\]

Podstawiając to, otrzymujemy:

\begin{align*}
P\!\left( \bigcap_{i \in I} (\Omega \setminus A_i) \right)
&= 1 - \left[ 1 - \prod_{i \in I} \big( 1 - P(A_i) \big) \right] \\
&= \prod_{i \in I} \big( 1 - P(A_i) \big) \\
&= \prod_{i \in I} P(\Omega \setminus A_i).
\end{align*}


\subsection*{Odpowiedź}
\[
\hfill \square
\]
